\begin{proposition}[C.3.18 (Proved): an actual bounding four-manifold and its frame defect]
There is a smooth genus-one Lefschetz fibration
\[
 \mathfrak g:Q\simeq S^5\longrightarrow S^3
\]
whose critical set is a circle mapping bijectively to the positive
$(3,4)$ knot
\[
 K=\{z_2^3-z_1^4=0\}\subset S^3.
\]
It agrees on the exterior and the outer cusp collar with the
actual torus quotient of C.3.13. In particular, the inverse image
of a nearby positive weighted circle $K'$ is the actual marked
three-torus with periods $(\gamma,u,w)$.

This three-torus bounds an embedded compact four-manifold
$M\subset Q$ carrying a positive genus-one Lefschetz fibration
over a Seifert surface of $K'$, with exactly twelve nodal critical
points. Its Euler number is $12$. Give $M$ its Lefschetz complex
orientation and its boundary the induced orientation, which in
the marked period convention is $(\gamma,w,u)$. Let
$\mathfrak f_{\mathrm{aff}}$ be the stable boundary frame consisting
of the three period directions and the two real $q$-normal
directions, with a fixed ordering consistent with the ambient
orientation. Relative to this actual affine frame,
\[
 \left\langle p_1(TM\oplus\R,\mathfrak f_{\mathrm{aff}}),
                   [M,\partial M]\right\rangle=-24.
\]
Equivalently, its transgression is $+24$ when the boundary
three-torus is oriented by the ordered basis $(\gamma,u,w)$.
No value of $\langle p_1(M_\psi)x,[M_\psi]\rangle$ or of
$J_\psi$ is inferred without the cap transgression calculation.
\end{proposition}
\begin{proof}
On the complement of the cusp body, the descended map of C.3.13
is a smooth two-torus bundle: on an ordinary chart it is the
quotient by the two residual real period directions $u,w$,
and the finite root charts give the same submersion after their
free finite quotients. The explicit finite gauges in C.3.13
ensure that this includes the two exceptional weighted circles
in the target.

Use the marked product of C.3.17 on the cusp body. In its
zero-unit model the projection is an elliptic $I_1$ neighbourhood
times a circle. Choose the product circle coordinate by the
second Tate identity section $z_2=1$ and the extending circle
action. On the boundary this coordinate is exactly the marked
character $\theta$: the section has value zero and the boundary
action is the actual $\gamma$ translation. The prescribed period
transport in C.3.17 preserves all these coordinates on the
outer collar. Consequently the map
\[
 (x,\theta)\longmapsto
 \Phi_0\bigl(q_E(x),\theta-\operatorname{Re}\kappa_0(q_E(x))\bigr)
\]
agrees there with the exterior map, where $q_E:E_{I_1}\to\Delta$
is the nodal elliptic map and $\Phi_0,\kappa_0$ are the exact
solid-torus coordinates of C.3.13. The two maps glue smoothly.
Near the critical circle, after the indicated smooth target
coordinate change, their local expression is
\[
 (x,y,\theta)\longmapsto(xy,\theta).
\]
Thus there is one positive nodal degeneration over each point
of $K$, and no other critical point. This constructs the stated
map from the actual marked collar. It need not equal the old
map $\overline f$ in the interior of the cusp body.

Choose $0<\epsilon\ll1$ and put
\[
 K'=\{(1-\epsilon)z_2^3-z_1^4=0\}\subset S^3,
 \qquad F_\epsilon=(1-\epsilon)z_2^3-z_1^4.
\]
Then $K'$ is the positive weighted circle $q=\epsilon$ in
the actual collar coordinates. The weighted circle field
$(z_1,z_2)\mapsto(e^{6\pi it}z_1,e^{8\pi it}z_2)$ satisfies
\[
 F_\epsilon(t\cdot z)=e^{24\pi it}F_\epsilon(z).
\]
Hence $\arg F_\epsilon$ is a submersion off $K'$. The closure
of one of its pages is a smooth compact oriented Seifert
surface $\Sigma'$ with boundary $K'$: the polynomial has a
simple complex normal zero along the link, which gives the
usual smooth half-plane collar of a page. Its fibres are
connected, since the normal meridian has argument degree one.
On the critical knot $K$,
\[
 F_\epsilon=-\epsilon z_2^3.
\]
Its argument increases twelve times along the positive weighted
parameter. Every page therefore meets $K$ transversely in
exactly twelve points, all with positive local intersection.

Set $M=\mathfrak g^{-1}(\Sigma')$. The displayed local nodal
model and transversality of $\Sigma'$ to the critical circle
show that $M$ is a smooth embedded compact four-manifold with
boundary $\mathfrak g^{-1}(K')$. The restriction to $M$ is a
positive genus-one Lefschetz fibration over $\Sigma'$ with
twelve critical points. A smooth torus has Euler number zero,
and an $I_1$ fibre has Euler number one. Stratification over
these twelve values therefore gives $\chi(M)=12$.

We next specify the boundary framing whose characteristic number
is computed. Choose a fibrewise complex structure on the regular
two-tori, compatible with their affine markings, and the usual
complex structure in the nodal charts. Such choices glue: on
the intervening collars the marked elliptic upper half-plane
is contractible. Choose a complex structure on $\Sigma'$ and
a compatible horizontal distribution. These give an almost
complex structure on $M$, equal to the standard nodal structure
near its critical points. Along $\partial M$, its vertical
complex tangent is framed by the translation vector dual to
the elliptic one-form with $w$-period one. Frame its horizontal
complex line by a lift of the outward normal of $\Sigma'$;
the associated real imaginary direction is the positive
boundary tangent. Denote this complex frame by $\mathfrak f_C$.

For completeness, the relative determinant class is pulled
back from the base. On a regular fibre its nonzero holomorphic
one-forms form a line $\lambda$, and evaluation identifies
the relative canonical line with the pullback of $\lambda$.
This remains true through each nodal chart: for $t=xy$ the
relative form is $dx/x=-dy/y$, and wedging with $dt$ gives
a nowhere-zero multiple of $dx\wedge dy$. Thus, as complex
line bundles,
\[
 \det(TM,J)\simeq\pi^*(T\Sigma'\otimes\lambda^{-1}).
\]
The determinant trivialization supplied by $\mathfrak f_C$
is the pullback of the corresponding base trivialization along
$\partial\Sigma'$. Therefore for some
$\beta\in H^2(\Sigma',\partial\Sigma';\Z)$,
\[
 c_1(TM,\mathfrak f_C)=\pi^*\beta,
 \qquad c_1(TM,\mathfrak f_C)^2=0.
\]
This does not require assigning a degree to $\lambda$ or a
holomorphic structure to the entire four-manifold. Its local
evaluation argument concerns the fibrewise elliptic structures
and the positive nodal charts.

The horizontal vector in $\mathfrak f_C$ points outwards on
the boundary. The relative top Chern class is therefore the
relative Euler class for an outward-pointing vector, and
Poincar\'e--Hopf gives
\[
 \langle c_2(TM,\mathfrak f_C),[M,\partial M]\rangle
 =\chi(M)=12.
\]
The relative characteristic-class identity now gives
\[
 \langle p_1(TM,\mathfrak f_C),[M,\partial M]\rangle
 =\langle c_1^2-2c_2,[M,\partial M]\rangle=-24.
\]

We finally compare this frame with the actual annular one.
On the target collar the exact formulas of C.3.13 give
\[
 F_\epsilon=(q-\epsilon)S(q)^{12}
              \exp(24\pi i\vartheta),
 \qquad\vartheta=\theta-\operatorname{Re}\kappa_0(q).
\]
Along $q=\epsilon$, the inward normal ray of the chosen page
therefore has $q$-angle equal to a constant minus $12\theta$.
Its complementary normal line in $S^3$ rotates by the same
angle. Add the pullback of that normal real line to $TM$,
identifying $TM\oplus\R$ with $TQ|_M$ after a horizontal
choice. On the boundary, the complex vertical translation
frame and the positive $\gamma$ tangent are the actual marked
period frame. The two remaining real vectors differ from
the fixed real and imaginary $q$ directions by a rotation
of winding $-12$ as $\gamma$ traverses the circle. Changes
of horizontal lifts add only contractible shear matrices.
Thus the stable frame comparison, up to fixed ordering, is
\[
 \operatorname{diag}\bigl(I_3,R_{-24\pi\theta}\bigr):
 T^3\longrightarrow SO(5).
\]
It factors through the $\gamma$ circle and is nullhomotopic,
since $\pi_1(SO(5))=\Z/2$ and twelve is even. Consequently
the relative first Pontryagin number is unchanged when
$\mathfrak f_C\oplus1$ is replaced by
$\mathfrak f_{\mathrm{aff}}$.

For the orientation convention, the elliptic complex orientation
is $(w,u)$, because its marked periods are $1,\tau$ with
$\operatorname{Im}\tau>0$. The outward-first boundary convention
for the Lefschetz orientation therefore orients the boundary
by $(\gamma,w,u)$. Reversing this to $(\gamma,u,w)$ reverses
the transgression sign, giving the equivalent $+24$ convention
stated above.
\end{proof}

\begin{remark}[Scope]
The four-manifold leaves the cusp body; its existence is
consistent with the local nonbounding assertion in C.3.17.
The value computed here is a relative Pontryagin defect of
the specified affine frame, not yet the Pontryagin pairing
of the closed cap in C.3.14. That latter conversion must
retain the actual gluing and its line-bundle transition.
\end{remark}
